% 2012 年陕西卷文科第 19 题：甲、乙两品牌产品使用寿命的频率分布直方图。
% 来源 PDF：普通高考/2012/2012陕西文.pdf，第 2 页（印面 766）第 2 栏。
% 重绘依据：甲品牌 [100,150)5、[150,200)20、[200,250)40、[250,300)25、[300,350]10；
%   乙品牌 [100,150)10、[150,200)20、[200,250)30、[250,300)25、[300,350]15。
\begin{tikzpicture}[scale=0.62, semithick, every node/.style={font=\small}]
  % ---------- 甲 ----------
  \begin{scope}
    \draw[->] (0,0) -- (6.6,0) node[below] {寿命/小时};
    \draw[->] (0,0) -- (0,4.8) node[left] {频率/组距};
    \foreach \x/\h in {1/0.5, 2/2.0, 3/4.0, 4/2.5, 5/1.0} {
      \draw[fill=gray!15] (\x,0) rectangle ++(1,\h);
    }
    \foreach \y/\t in {0.5/5,1/10,2/20,2.5/25,4/40} {
      \draw (0,\y) -- (-0.08,\y) node[left] {\t};
    }
    \foreach \x/\t in {1/100,2/150,3/200,4/250,5/300,6/350} {
      \draw (\x,0) -- (\x,-0.08) node[below] {\t};
    }
    \node at (3,-1.1) {甲品牌};
  \end{scope}
  % ---------- 乙 ----------
  \begin{scope}[shift={(9.0,0)}]
    \draw[->] (0,0) -- (6.6,0) node[below] {寿命/小时};
    \draw[->] (0,0) -- (0,4.8) node[left] {频率/组距};
    \foreach \x/\h in {1/1.0, 2/2.0, 3/3.0, 4/2.5, 5/1.5} {
      \draw[fill=gray!15] (\x,0) rectangle ++(1,\h);
    }
    \foreach \y/\t in {1/10,1.5/15,2/20,2.5/25,3/30} {
      \draw (0,\y) -- (-0.08,\y) node[left] {\t};
    }
    \foreach \x/\t in {1/100,2/150,3/200,4/250,5/300,6/350} {
      \draw (\x,0) -- (\x,-0.08) node[below] {\t};
    }
    \node at (3,-1.1) {乙品牌};
  \end{scope}
\end{tikzpicture}
